% =====================================================================
%   附录 H：网页购物 (Web Shopping)
% =====================================================================

\section{网页购物}\label{app:ws}

\subsection{数据集细节}

\subsubsection*{构造细节}

环境向智能体展示网页的文本观测及可选动作。智能体可如同在真实世界中一样
自由探索网站，并通过可点击的按钮浏览商品。约一百万件商品从 \code{amazon.com}
爬取，构成网站的数据库。随后，每件商品均被打上反映其自身属性的标签。
共收集了 12\,087 条人类指令，并将每条指令与目标及期望属性相关联。
数据集构造细节请参见 \citet{yao2022webshop}。

\subsubsection*{评测设置}

我们采用 12\,087 条指令中的前 500 条作为测试集（沿用
\citet{yao2022webshop}~官方实现）。每一轮交互可分解为如下步骤：
\begin{enumerate}
    \item \textbf{指令下达}：在初始提示之后（提示中包含环境信息以及 LLM
          应采用的回复格式），我们给出关于期望购买的商品的指令。
    \item \textbf{交互}：智能体按提示以规定的格式进行响应，包含其思考
          与拟采取的动作。可选动作分为两类：\code{search} 与 \code{click}，
          分别对应于现实中的搜索与点击按钮。环境以简化版的网页文本与
          可用按钮列表作为响应。该过程将反复进行，直至智能体点击
          \code{buy now} 按钮或超出轮数限制。
    \item \textbf{奖励计算}：我们采用原论文中的奖励函数作为评测指标。
          该奖励将期望属性与被购商品实际属性之间的相似度映射为 0 到 1
          之间的数值。
\end{enumerate}

\subsubsection*{评测指标}

鉴于给定查询可能对应多个合适的商品，Webshop 采用一种匹配奖励作为评测指标：
\begin{equation}\label{eq:webshop-reward}
    \text{Reward} = \frac{|U_{\text{att}} \cap Y_{\text{att}}| + |U_{\text{opt}} \cap Y_{\text{opt}}| + \mathbb{I}[y_{\text{price}} \le u_{\text{price}}]}
                          {|U_{\text{att}}| + |U_{\text{opt}}| + 1}
                  \cdot r_{\text{type}},
\end{equation}
其中
\begin{equation}\label{eq:webshop-rtype}
    r_{\text{type}} = \begin{cases}
        0,   & \text{若 TextMatch} = 0, \\
        0.1, & \text{若 TextMatch} < 0.1, \\
        0.5, & \text{若 TextMatch} \le 0.2\ \text{且 query 不匹配且 category 不匹配}, \\
        1,   & \text{其它情形}.
    \end{cases}
\end{equation}
式中，$U$ 与 $Y$ 分别表示目标商品与所购商品的属性集；att 与 opt 分别指
属性与选项；TextMatch 为所购商品标题与目标商品标题之间的代词、名词与专有名词
匹配度。

\subsection{提示词示例}

我们使用如下格式的提示词：
\begin{quote}\small\ttfamily
User:\\
You are web shopping. I will give you instructions about what to do.
You have to follow the instructions. Every round I will give you an
observation and a list of available actions, you have to respond an
action based on the state and instruction. You can use search action if
search is available. You can click one of the buttons in clickables.\\
\\
An action should be of the following structure:\\
\code{search[keywords]}\\
\code{click[value]}\\
\\
If the action is not valid, perform nothing. Keywords in search are up
to you, but the value in click must be a value in the list of available
actions. Remember that your keywords in search should be carefully
designed.\\
\\
Your response should use the following format:\\
Thought:\\
I think \dots\\
Action:\\
\code{click[something]}\\
\\
User:\\
Observation: \code{\{observation\}}\\
\\
Available Actions: \code{\{available_actions\}}\\
\end{quote}

该数据集采用 1-shot 示例（详见原论文）。下面是若干用户—智能体交互示例的摘要：
\begin{itemize}
    \item 初始 Observation 包含指示文本、可用动作的 JSON；智能体决定
          \code{search[...]}
    \item 后续 Observation 给出商品列表；智能体选择目标商品的编号进行
          \code{click[...]}
    \item 进入商品详情页后，智能体选择符合需求的规格（如尺寸），再次
          \code{click[...]}
    \item 最终智能体通过 \code{click[Buy Now]} 完成购买。
\end{itemize}
完整对话示例及对应英文原版提示词请参见原论文附录 H。
